\chapter{A Universe of Pressure and Geometry: The SDT Cosmology}

\author{James Tyndall}

%----------------------------------------------------------------------
\section{Introduction: The Cosmos as a Unified Mechanical System}
%----------------------------------------------------------------------

The prevailing $\Lambda$CDM model relies on extraordinary, unproven postulates: a singular Big Bang from nothing, superluminal inflation, and a universe dominated by invisible dark matter and dark energy.

SDT offers a fundamentally different vision: the universe is an \textbf{eternal, infinite, and mechanically evolving system}.  Its observed properties are the current state of a continuous, universal process of matter and pressure finding geometric equilibrium.


%----------------------------------------------------------------------
\section{Galactic Structure: A Dynamic Pressure Vortex}
%----------------------------------------------------------------------

A galaxy is not a static collection of stars held by an invisible dark matter halo.  It is a \textbf{rotating displacement vortex} governed by the dynamic, relativistic effects of its own propagating pressure field.

\subsection{The Flat Rotation Curve}

\textbf{Mechanism:}
\begin{enumerate}
  \item The galaxy's rotating central mass creates a \textbf{spiral pressure wave} propagating outwards at $c$.
  \item This spiral wave exerts a continuous, non-zero \textbf{tangential force} on outer stars.
  \item This ``push'' transfers angular momentum from inner to outer regions.
\end{enumerate}

\textbf{Consequence:} The continuous energy input prevents Keplerian velocity decay.  The ``missing mass'' of dark matter is an illusion from applying a static force law to a dynamic, relativistic system.

\begin{equation}\label{eq:v-galaxy}
  v_{\text{SDT}}(r) = \sqrt{r \cdot a_r(r)}
\end{equation}

where $a_r$ is the radial acceleration from the full, retarded potential of the galaxy's baryonic mass.


%----------------------------------------------------------------------
\section{The Nature of Cosmological Redshift}
%----------------------------------------------------------------------

Redshift is not primarily a Doppler effect from metric expansion.  It is a physical consequence of \textbf{energy dilution}.

\textbf{Mechanism:} A photon (percussive pressure pulse) expands as a spherical wavefront.  Total energy is conserved, but it is spread over an increasing surface area:
\begin{equation}\label{eq:dilution}
  \varepsilon(r) = \frac{E_{\text{total}}}{4\pi r^2}
\end{equation}

The energy a telescope intercepts, and the frequency it observes ($E = hf$), is a function of this diluted energy density.  The non-linear distance-redshift relationship naturally appears as ``accelerating expansion'' when misinterpreted through a Doppler lens.


%----------------------------------------------------------------------
\section{The Cosmic Microwave Background}
%----------------------------------------------------------------------

The CMB is not the ``afterglow of the Big Bang.''  It is the \textbf{observable horizon of our infinite universe.}

\textbf{Mechanism:} The CMB is the collected light from the most distant galaxies whose pressure waves have been diluted into the microwave spectrum.

\textbf{Consequence:} The CMB is not a ``surface of last scattering'' in time; it is a \textbf{sphere of current observation in space}.  Its isotropy is the ultimate proof of the cosmological principle: in an infinite, eternal universe, the view from any point must, on average, look the same in every direction.  The tiny anisotropies are the imprints of the most ancient, large-scale pressure occlusions in the spation medium.


%----------------------------------------------------------------------
\section{Black Holes and the Cyclical Universe}
%----------------------------------------------------------------------

SDT rejects the concept of a singularity.

\subsection{Black Holes as ``Darkstars''}

A black hole is a \textbf{Darkstar}---a celestial object where matter has reached maximum displacement density.  It is a perfectly packed, crystalline arrangement of neutron vortices.  The event horizon is where the propagated pressure field creates a local escape velocity of $c$.  Information is not destroyed; it is \textbf{repacked} into maximum geometric order.

\subsection{The ``Big Bang'' as a Phase Transition}

Darkstars are not eternal.  They accumulate stress from background pressure.  At a critical super-massive state, they undergo a \textbf{cosmic phase transition}:
\begin{itemize}
  \item Internal geometric constraints shift, causing catastrophic loss of stability.
  \item The Darkstar ``depressurises'' explosively, releasing perfectly ordered matter as a plasma of fundamental vortices.
\end{itemize}

This event---appearing to an observer as a ``Big Bang''---is not creation \emph{ex nihilo}, but a local, cyclical act of \textbf{repacking and renewal}.  The universe is an eternal, steady-state system punctuated by local, violent recycling events.


%----------------------------------------------------------------------
\section{Conclusion: A Universe of Geometry and Pressure}
%----------------------------------------------------------------------

SDT cosmology describes a universe that is:
\begin{itemize}
  \item \textbf{Mechanical, not Abstract:} governed by tangible pressure and displacement.
  \item \textbf{Deterministic, not Probabilistic:} a causal chain from particle to cosmos.
  \item \textbf{Eternal and Cyclical, not Singular:} without beginning or end, in constant evolution.
\end{itemize}

By providing mechanical explanations for galactic rotation, cosmological redshift, the CMB, and the nature of black holes, SDT eliminates the need for dark matter, dark energy, and a singular Big Bang.


%----------------------------------------------------------------------
\section*{And Therefore}
%----------------------------------------------------------------------

\textbf{Therefore\ldots} the universe is comprehensible.  The paradoxes of modern physics are artifacts of an incomplete framework.

\textbf{Therefore\ldots} the fundamental forces are unified.  They are a tiered hierarchy stemming from a single interaction: the displacement of space by matter.

\textbf{Therefore\ldots} the quantum and the cosmic are one.  The same laws that govern the stability of an atom govern the rotation of a galaxy.

\textbf{Therefore\ldots} the universe is not made of ``things,'' but of \textbf{patterns}---an architecture of stable, resonant geometries in a single, dynamic medium.

\textbf{And therefore\ldots} a new era of physics is possible.  An era based on the deterministic, tangible, and elegant principles of geometry and motion.

The work is not finished.  It has just begun.  We have laid the foundation.  Now, we begin to build the future.
